\documentclass[11pt]{article}
\usepackage[margin=1in]{geometry}
\usepackage{amsmath,amssymb,amsthm,hyperref}
\newtheorem{theorem}{Theorem}
\newtheorem{lemma}{Lemma}
\newtheorem{corollary}{Corollary}
\newcommand{\E}{\mathbb E}
\title{Positive coefficients in a centered exponential remainder\\
MathDB 364606}
\author{Research record in Principia Math in Progress}
\date{12 September 2026}
\begin{document}
\maketitle

\section*{Task 1: solve the nonnegative coefficient conjecture}
For an integer $r\geq1$, set
\[
 B_r(x)=e^x-\sum_{j=0}^{r-1}\frac{x^j}{j!}.
\]
Conjecture 3.2 of Canfield, Helton and Hughes~\cite{CHH}, on published
page 241, asserts that the Maclaurin coefficients of
$e^{-x/(r+1)}B_r(x)$ are all nonnegative.

\begin{theorem}\label{thm:main}
The conjecture is true. For $r>1$, the coefficients below degree $r$
and at degree $r+1$ vanish, the coefficient at degree $r$ is $1/r!$,
and every coefficient of degree at least $r+2$ is strictly positive.
For $r=1$,
\[
 e^{-x/2}B_1(x)=2\sinh(x/2),
\]
so precisely the positive odd degrees have positive coefficients.
Moreover, for any real $c$,
\[
 e^{-cx}B_r(x)\ \hbox{has only nonnegative Maclaurin coefficients}
 \quad\Longleftrightarrow\quad c\leq\frac1{r+1}.
\]
\end{theorem}

\paragraph{Relation to existing work.}
The proof identifies the coefficients with moments of a centered
$\mathrm{Beta}(1,r)$ variable. Positivity of odd beta central moments
when the second shape parameter exceeds the first is already known:
see Skorski~\cite[Theorem 3 and Corollary 3]{Skorski}.
Beta-moment representations for associated Stirling numbers also appear
in Gismatullin and Tardivel~\cite{GT}. We claim no novelty for those
probabilistic tools; the argument applies them to this centered
coefficient conjecture.
For completeness, the needed special case is proved below directly by
comparing the density on the two sides of its mean.

\begin{lemma}\label{lem:moments}
Let $T$ have density $f(t)=r(1-t)^{r-1}$ on $[0,1]$, where $r\geq1$,
and let $\mu=\E T=1/(r+1)$. Then
$m_n=\E(T-\mu)^n$ satisfies $m_0=1$, $m_1=0$, and
$m_{2k}>0$ for $k\geq1$.
For $r>1$, also $m_{2k+1}>0$ for every $k\geq1$;
for $r=1$, every odd moment vanishes.
\end{lemma}

\begin{proof}
The density integrates to one and direct integration gives
$\E T=1/(r+1)$. The even moment assertions follow from nondegeneracy.
If $r=1$, reflection about $1/2$ proves the odd moment assertion.

Suppose $r>1$, so that $\mu<1/2$ and $f$ is strictly decreasing on
$(0,1)$. For $0<s<1-\mu$ define
\[
 q(s)=
 \begin{cases}
 f(\mu+s)-f(\mu-s),&0<s<\mu,\\
 f(\mu+s),&\mu<s<1-\mu .
 \end{cases}
\]
Values at the endpoints do not affect the integrals. Splitting at the
mean and changing variables gives, for every odd positive integer $n$,
\[
 m_n=\int_0^{1-\mu}s^nq(s)\,ds.
\]
In particular $\int_0^{1-\mu}s q(s)\,ds=m_1=0$.
We have $q(s)<0$ below $\mu$ and $q(s)>0$ above $\mu$.
Consequently, for $k\geq1$,
\[
 m_{2k+1}
 =\int_0^{1-\mu}s\bigl(s^{2k}-\mu^{2k}\bigr)q(s)\,ds>0.
\]
The integrand is positive on both open intervals $(0,\mu)$ and
$(\mu,1-\mu)$, which have positive length.
\end{proof}

\begin{proof}[Proof of Theorem \ref{thm:main}]
The integral remainder formula for the exponential gives, for every
complex $x$,
\[
 B_r(x)=\frac{x^r}{(r-1)!}
        \int_0^1(1-t)^{r-1}e^{xt}\,dt
       =\frac{x^r}{r!}\E e^{xT}.
\]
The equality can also be checked term by term using
$\E T^n=r!\,n!/(r+n)!$.
Since $T$ is bounded, the exponential series converges uniformly for
$x$ in compact subsets of $\mathbb C$ and $t\in[0,1]$. Thus
\begin{equation}\label{eq:coefficients}
 e^{-\mu x}B_r(x)=
 \frac{x^r}{r!}\E e^{x(T-\mu)}
 =\sum_{n=0}^{\infty}\frac{m_n}{r!\,n!}x^{r+n}.
\end{equation}
Lemma \ref{lem:moments} proves all the asserted signs and zeros.
The expression for $r=1$ follows directly from $B_1(x)=e^x-1$.

If $c\leq\mu$, write
\[
 e^{-cx}B_r(x)=e^{(\mu-c)x}
              \bigl(e^{-\mu x}B_r(x)\bigr).
\]
Both factors have nonnegative coefficients. Conversely, the coefficient
of $x^{r+1}$ in this product is $(\mu-c)/r!$, which is negative when
$c>\mu$. This proves the equivalence.
\end{proof}

\paragraph{An independent recurrence check.}
There is a second short way to obtain the moment signs. Put
$w(t)=t(1-t)f(t)=rt(1-t)^r$. Its endpoint values vanish and
$w'(t)=[1-(r+1)t]f(t)$. Integration by parts with
$g(t)=(t-\mu)^{n-1}$, for $n\geq2$, gives
\[
 (r+1)m_n
 =(n-1)\{\mu(1-\mu)m_{n-2}+(1-2\mu)m_{n-1}-m_n\},
\]
and hence
\begin{equation}\label{eq:recurrence}
 (r+n)m_n=(n-1)
 \left\{\frac{r}{(r+1)^2}m_{n-2}
       +\frac{r-1}{r+1}m_{n-1}\right\}.
\end{equation}
This is the $\mathrm{Beta}(1,r)$ specialization of the known recurrence
in~\cite[Theorem 3]{Skorski}. Its nonnegative coefficients and the initial
values $m_0=1,m_1=0$ supply another all-order proof.
The accompanying exact-arithmetic script compares this recurrence,
the binomial expansion of the moments, and the original product series.
Those finite checks supplement the analytic proof.

\section*{Follow-up tasks after the full solution}
\paragraph{Task 2.}
Obtain two-sided explicit bounds for
$[x^{r+n}]e^{-x/(r+1)}B_r(x)=m_n/(r!\,n!)$ that are uniform in
integers $r\geq2$ and $n\geq2$, with asymptotically sharp constants
in both the fixed-$r$ and fixed-$n$ regimes.

\paragraph{Task 3.}
Find a uniform asymptotic expansion of $m_n$ when $r,n\to\infty$
and $n/r$ remains in a compact subinterval of $(0,\infty)$.
Identify the different contributions of the two sides of the mean,
including any difference between even and odd $n$.

\paragraph{Task 4.}
Use the coefficient positivity proved here to simplify the estimate
$|J|\to0$ singled out in Section 3.3.1 of~\cite{CHH}, and determine
whether this yields an explicit uniform error estimate for the
associated Stirling-number approximation.

\paragraph{Task 5.}
For $r\geq2$ and $u\geq0$, let $T_{r,u}$ have density proportional to
$e^{ut}(1-t)^{r-1}$ on $[0,1]$. Determine exactly which pairs $(r,u)$
make every coefficient of
$\E\exp\{x(T_{r,u}-\E T_{r,u})\}$ nonnegative.

\paragraph{Status and provenance.}
The February 2026 published source retains the target as a conjecture.
On 12 September 2026 the checked MathDB record was open with no claims
or discussion, and searches of the current primary literature and
public claim catalogues located no full claim for this coefficient
conjecture. The known beta-moment theorem is explicitly acknowledged.
The accompanying dated audit records the scope and limits of that
search; no exhaustive novelty guarantee is intended.

\begin{thebibliography}{9}
\bibitem{CHH}
E. R. Canfield, J. W. Helton and J. A. Hughes,
\emph{Uniform convergence of an asymptotic approximation to associated
Stirling numbers}, Contributions to Discrete Mathematics
\textbf{21}(1) (2026), 227--244.
\href{https://doi.org/10.55016/dptxws28}{doi:10.55016/dptxws28}.
Preprint: \href{https://arxiv.org/abs/2409.01489}{arXiv:2409.01489}.
\bibitem{Skorski}
M. Skorski, \emph{Bernstein-type bounds for beta distribution},
Modern Stochastics: Theory and Applications \textbf{10}(2) (2023),
211--228.
\href{https://doi.org/10.15559/23-VMSTA223}{doi:10.15559/23-VMSTA223}.
\bibitem{GT}
J. Gismatullin and P. Tardivel,
\emph{Beta distribution and associated Stirling numbers of the second kind},
Probability and Mathematical Statistics \textbf{44}(1) (2024), 119--132.
\href{https://doi.org/10.37190/0208-4147.00156}{doi:10.37190/0208-4147.00156}.
\end{thebibliography}
\end{document}
